A spacecraft is orbiting the Near Earth Asteroid (2608) \textit{Seneca}
(staying continuously very close to the asteroid), transmitting pulsed data to
the Earth. Due to the relative motion of the two bodies (the asteroid and the
Earth) around the Sun, the time it takes for a pulse to arrive at the ground
station varies approximately between 2 and 39 minutes. The orbits of the Earth
and Seneca are coplanar. Assuming that the Earth moves around the Sun on a
circular orbit (with radius $a_{\mathrm{Earth}} = 1$\,AU and period
$T_{\mathrm{Earth}} = 1$\,yr) and that the orbit of Seneca does not intersect
the orbit of the Earth, calculate:

\begin{description}
\item[(1)] the semi-major axis, $a_{\mathrm{Sen}}$, and the eccentricity,
  $e_{\mathrm{Sen}}$, of Seneca's orbit around the Sun;
\item[(2)] the period of Seneca's orbit, $T_{\mathrm{Sen}}$, and the average
  period between two consecutive oppositions, $T_{\mathrm{syn}}$, of the
  Earth--Seneca couple;
\item[(3)] an approximate value for the mass of the planet Jupiter,
  $M_{\mathrm{Jup}}$ (assuming this is the only planet of our Solar system with
  non-negligible mass compared to the Sun). Assume that the presence of
  Jupiter does not influence the orbit of Seneca.
\end{description}
